\documentclass[11pt]{article}

\usepackage[utf8]{inputenc}
\usepackage[T1]{fontenc}
\usepackage{graphicx} % Required for inserting images
\usepackage[a4paper,margin=1in]{geometry}
\usepackage{amsmath}
\usepackage{amssymb}
\usepackage{array}
\usepackage{tabularx}
\usepackage{pdflscape}
\usepackage{needspace}
\usepackage[hidelinks]{hyperref}
\emergencystretch=3em % absorb overlong lines from long citation keys

\title{project-36}
\author{ELITE Research Lab LLC}
\date{July 2026}

\begin{document}
	
	\maketitle
	
	\section{Review Methodology and the Secondary Evidence Base}
	
	Ten of the sixty-six studies in this corpus (15.2\%) are reviews, surveys, or conceptual syntheses. They constitute evidence about the field's methodology rather than instances of it, and they supply the vocabulary in which the remaining fifty-six studies are described. Seven systematic or scoping reviews of large language models (LLMs) in mental health appeared between 2024 and 2026, so a further synthesis must show what those seven leave unaddressed. The ten papers fall into three groups: scoping reviews that map breadth, systematic reviews conducted under PRISMA that also attempt appraisal, and conceptual papers that supply taxonomies and governance frameworks without primary data.
	
	\subsection{Scoping reviews}
	
	Four studies adopt a scoping design aimed at coverage rather than effect estimation.
	
	Jin et al. (2025)\cite{Jin2025LLMMentalHealth} searched seven databases, including the Chinese-language Weipu, CNKI and Wanfang, for work published between 1 January 2019 and 31 August 2024, including 95 of 4,859 retrieved records. Two features distinguish it: prospective registration on the Open Science Framework, adopted by only two of the seven empirical reviews, and systematic non-English coverage, which matters given the volume of primary work from Chinese institutions. Screening proceeded in three rounds with disagreements escalated to third reviewers, and extracted data were synthesised descriptively in R 4.4.2 without pooling.
	
	Hua et al. (2025a)\cite{Hua2025LLMMentalHealthScoping} searched six sources, three of them preprint servers, covering 1 October 2019 to 2 December 2023, and included 34 of 313 records. Admitting preprints reduces the recency lag affecting every other review here at the cost of unrefereed work. Its most consequential contribution is internal to its own process: GPT-4 served as a parallel screening reviewer alongside a human, with agreement quantified by Cohen's $\kappa$ ($\approx$ 0.90). This is the only study in the corpus to treat LLM-assisted evidence synthesis as an object of measurement rather than a convenience.
	
	Yang et al. (2026)\cite{fpsyg.2025.1715306} retrieved 10,743 records across eleven databases, assessed 58 at full text and charted 29, following both PRISMA-ScR and the Joanna Briggs Institute framework. Two independent reviewers screened and charted, with a third resolving discrepancies. Where the other reviews ask what LLMs can do, this one asks where their application boundaries lie, charting technical, ethical and practical limits per study.
	
	Hua et al. (2025b)\cite{Hua2025} searched four databases for peer-reviewed work published between 1 January 2020 and 19 July 2024, including 16 of 726 unique articles under PRISMA 2020. Its narrower scope, generative tasks only under a parameter threshold, yields the smallest included set among the scoping reviews. Mapping primary studies onto a three-level clinical evaluation pyramid converts a descriptive exercise into a diagnostic one and shows that the literature clusters at the lowest tier.
	
	\subsection{PRISMA systematic reviews}
	
	Three studies adopt a systematic design, and the differences between them are instructive precisely because the label is shared.
	
	Wang et al. (2025a)\cite{WANG2025} screened 1,046 references across five databases with paired independent reviewers in Rayyan and included eight studies, the most restrictive inclusion rate in the corpus at 0.76\%. It is the only review to apply a formal quality appraisal instrument, the Mixed Methods Appraisal Tool. Classifying primary research by the clinical competencies a human therapist would require reframes model evaluation around clinical role rather than technical task. With $n=8$, and with included studies weighted toward zero-shot prompt tests, its conclusions rest on a thin base.
	
	Kolding et al. (2025)\cite{Kolding2025} screened 724 unique records after removing 432 duplicates from 1,156, including 40 studies under a protocol preregistered on the Open Science Framework, with dual-reviewer screening at both stages supported by Covidence. The authors state that no formal quality appraisal or quantitative synthesis was performed, attributing this to the immaturity of the field. The candour is useful: it identifies the field's condition rather than the reviewers' resources as the binding constraint. Their finding that models and prompts are frequently underspecified to the point where replication is impossible is corroborated by the extraction reported in Section 7.
	
	Guo et al. (2024a)\cite{GUO2024} collated 40 peer-reviewed articles from five databases across the widest temporal window in the corpus, 1 January 2017 to 30 April 2024. It is the only review to apply the Cochrane Risk of Bias 2 tool, and one of two to report inter-reviewer agreement (Cohen's $\kappa$ = 0.84). Its conclusion is the most restrictive here, that current risks of clinical use may outweigh the benefits, and it identifies the absence of a benchmarked ethical framework as a structural gap alongside hallucination and output inconsistency. Its acknowledged limitation is that 98\% of reviewed papers cover only BERT- and GPT-family models.
	
	\needspace{8\baselineskip}
	\begin{center}
		\textbf{Table 1.1: Review protocol parameters across the seven empirical reviews}
		\vspace{2pt}
		
		\tiny
		\setlength{\tabcolsep}{2pt}\renewcommand{\arraystretch}{1.1}
		\begin{tabularx}{\textwidth}{|X|X|X|X|X|X|X|X|X|}
			\hline
			\textbf{Study} & \textbf{Design} & \textbf{Databases} & \textbf{Screened $\rightarrow$ Included} & \textbf{Search window ends} & \textbf{Dual screening} & \textbf{Quality appraisal} & \textbf{Agreement reported} & \textbf{Pooling} \\
			\hline
			Jin et al. (2025)\cite{Jin2025LLMMentalHealth} & Scoping & 7 (incl. 3 Chinese) & 4,859 $\rightarrow$ 95 & Aug 2024 & Yes (3 rounds) & No & No & None \\
			\hline
			Hua et al. (2025a)\cite{Hua2025LLMMentalHealthScoping} & Scoping & 6 (incl. preprints) & 313 $\rightarrow$ 34 & Dec 2023 & Human + GPT-4 & No & $\kappa$ $\approx$ 0.90 & None \\
			\hline
			Yang et al. (2026)\cite{fpsyg.2025.1715306} & Scoping (PRISMA-ScR + JBI) & 11 & 10,743 $\rightarrow$ 29 & Not reported & Yes (2 + adjudicator) & No & No & None \\
			\hline
			Hua et al. (2025b)\cite{Hua2025} & Scoping (PRISMA 2020) & 4 & 726 $\rightarrow$ 16 & Jul 2024 & Yes & No & No & None \\
			\hline
			Wang et al. (2025a)\cite{WANG2025} & Systematic (PRISMA) & 5 & 1,046 $\rightarrow$ 8 & Jun 2024 & Yes (paired) & MMAT & No & None \\
			\hline
			Kolding et al. (2025)\cite{Kolding2025} & Systematic (PRISMA, preregistered) & 3 & 724 $\rightarrow$ 40 & Feb 2024 & Yes (both stages) & Declined, stated & No & None \\
			\hline
			Guo et al. (2024a)\cite{GUO2024} & Systematic (PRISMA) & 5 & Not reported $\rightarrow$ 40 & Apr 2024 & Yes (2 + adjudicator) & RoB 2 & $\kappa$ = 0.84 & None \\
			\hline
		\end{tabularx}
	\end{center}
	
	\subsection{Conceptual and governance contributions}
	
	Three papers contribute no primary data but supply instruments used throughout this chapter.
	
	Huang et al. (2025)\cite{Huang_2025} survey hallucination in LLMs. Although general-domain, it provides the definitional apparatus applied in Section 6: a layered taxonomy separating factuality hallucination (contradiction, fabrication) from faithfulness hallucination (instruction, context and logical inconsistency), with root causes mapped across data, training and inference stages and matched to detection and mitigation strategies. Residual hallucination surviving retrieval augmentation receives dedicated treatment, which bears directly on the retrieval-augmented systems examined in Section 3.
	
	Chung et al. (2023)\cite{chung2023challengeslargelanguagemodels} map five architectural challenges, hallucination, interpretability, bias, privacy and clinical effectiveness, to practical mitigations, drawing on deployment experience. Adapting the intrinsic/extrinsic hallucination distinction to counselling dialogue is the most useful move for present purposes: intrinsic hallucination contradicts the dialogue history, whereas extrinsic hallucination cannot be verified against any available input. The authors propose pairing stochastic neural generation with deterministic hard-coded overrides for crisis detection, conceding that hallucination cannot be eliminated within the current paradigm and must instead be contained.
	
	Asman et al. (2025)\cite{info:doi/10.2196/70439} frame a theme issue on responsible generative AI in mental health, offering a governance roadmap that integrates a relational ethics of care with clinical alliance standards, together with a proposed clinical risk assessment questionnaire for patients and developers.
	
	\subsection{Appraisal of the secondary evidence}
	
	Four properties recur across the seven empirical reviews and define the gap this work addresses. None performs quantitative synthesis, a defensible response to genuine heterogeneity given that the primary literature reports F1, balanced accuracy, AUC, BLEU, ROUGE, $\kappa$, ICC and bespoke Likert instruments without shared task definitions, but the consequence is that no pooled effect estimate exists for any LLM mental health application. Quality appraisal is the exception rather than the norm: two of seven apply a formal instrument and two report inter-reviewer agreement, so five synthesise studies of unassessed quality. Inclusion rates range from 0.27\% to 10.9\%, a fortyfold spread across overlapping periods and topics, which appears to reflect differing operational definitions of "LLM" and "mental health application" rather than differing literatures.
	
	Every search window closes before the majority of this corpus was published. The latest end-date among the seven is August 2024 and the earliest December 2023, whereas 45 of 66 studies here (68\%) appeared in 2025 or 2026 (Table 1.2). The existing secondary literature therefore cannot speak to multi-turn adversarial safety benchmarking, agentic red teaming, LLM-as-judge reliability frameworks, or preference-optimisation approaches to alignment.
	
	\needspace{8\baselineskip}
	\begin{center}
		\textbf{Table 1.2: Corpus publication years relative to prior review coverage}
		\vspace{2pt}
		
		\footnotesize
		\setlength{\tabcolsep}{3pt}\renewcommand{\arraystretch}{1.1}
		\begin{tabularx}{\textwidth}{|X|X|X|X|}
			\hline
			\textbf{Year} & \textbf{Studies in corpus} & \textbf{Share} & \textbf{Covered by any prior review's search window} \\
			\hline
			2023 & 4 & 6\% & Yes \\
			\hline
			2024 & 17 & 26\% & Partially \\
			\hline
			2025 & 36 & 55\% & No \\
			\hline
			2026 & 9 & 14\% & No \\
			\hline
			\textbf{Total} & \textbf{66} & \textbf{100\%} & \textbf{32\% at most} \\
			\hline
		\end{tabularx}
	\end{center}
	
	A fifth observation concerns scope. Five of seven reviews organise around application areas. Only Yang et al., through its boundary framing, and Hua et al. (2025b), through its evaluation pyramid, organise around evaluative adequacy. None organises around methodology as such: how corpora are constructed, how models are adapted, how evaluation is designed, and how risk is assessed. That is the organising principle adopted here.
	
	\subsection{Positioning of the present review}
	
	Five commitments follow from the appraisal above. Extraction is methodology-only, so that where prior reviews chart what models achieved, this review charts how the achievement was established. Every study is extracted across the same twenty-four fields, including four risk fields (hallucination, bias, safety, human evaluation) that no prior review extracts systematically, which permits the corpus-wide statements made in Sections 6 and 7. The search window extends through 2026, capturing the 68\% of the corpus outside all prior coverage. Absence is recorded as data, with "Not Reported" distinguished from "not applicable", so that gaps in the primary literature become findings. Two studies appeared in both preprint and published form and were merged into single records, giving 66 unique studies from 68 retrieved documents.
	
	Two limitations should be stated at the outset. Like all seven reviews examined here, this review performs no quantitative pooling, for the same heterogeneity reason. No formal quality appraisal instrument is applied to included studies; methodological quality is instead characterised through the extraction fields themselves and appraised cross-sectionally in Section 7. This is a weaker guarantee than a validated instrument would provide and represents a deliberate trade-off against extraction depth.
	
	\section{Data Foundations and Corpus Construction}
	
	Data constraints determine which methods are available. The recurring criticism that models are evaluated on social media proxies rather than clinical presentations is usually stated as a limitation, but read across twelve studies whose primary contribution is a corpus or an annotation protocol it is better understood as a structural condition. Clinical mental health data is scarce, legally encumbered and small, and most downstream choices examined in Sections 3 to 6 respond to that scarcity. Twelve studies (18.2\%) are grouped here because their principal contribution is the construction, integration or annotation of a dataset rather than a modelling technique.
	
	\subsection{Clinical and interview corpora}
	
	Four studies work with data generated inside a clinical or semi-clinical encounter, and with one exception they are the smallest datasets in the corpus.
	
	Ohse et al. (2024)\cite{OHSE2024101663} collected the KID corpus of German semi-clinical interviews, transcribed with Whisper large-v2. Fine-tuning on DAIC-WOZ and Extended DAIC makes the study simultaneously a corpus contribution and a generalisation test across language and protocol. Ground truth was PHQ-8 self-report, dichotomised at the conventional cut-off for classification and retained continuously for severity estimation. The authors note a unimodal text approach discarding prosodic and visual cues, translation-induced loss of nuance, and age and gender imbalance in recruitment.
	
	Sood et al. (2023)\cite{10385480} addressed scarcity by construction, integrating DAIC, E-DAIC and the Chinese-language EATD corpus into I-DAIC. Machine translation of EATD via the M2M 1.2B model was validated by bilingual annotators, making this the only cross-lingual merge here with a documented verification step. Five strategies for handling documents exceeding transformer context limits are benchmarked, from truncation variants to sliding windows and extractive summarisation. Fine-tuned BERT and RoBERTa underperformed TF-IDF-based SVM and logistic regression on weighted F1, attributed to class imbalance, which tempers any assumption that transformer adaptation is automatically superior on small clinical corpora.
	
	Xu et al. (2025a)\cite{Xu2025.01.03.24318117} assembled the largest genuinely clinical dataset here: psychiatrist-patient dialogues from 1,160 outpatients at the Shanghai Mental Health Center, roughly 15,000 minutes of transcribed Mandarin audio with electronic medical records covering 25 diagnoses. Ground truth derives from psychiatrist documentation rather than self-report, validated against 50 EMR cases reviewed by specialist psychiatrists. Methodologically it is the most complete study in the group, combining 138 validated clinical features and 177 scale items with an ensemble classifier over LLM outputs, LIWC counts, TF-IDF weights and OpenSMILE prosodic features, SMOTE oversampling, 10-fold cross-validation, Mann-Whitney U tests with FDR correction, and McNemar's test. Clinical noise degraded the acoustic features and Mandarin-centric training restricts transfer.
	
	Xin et al. (2024)\cite{xin2024outcomes} used IMPACT-ME, 102 end-of-treatment interviews from 34 sets of depressed adolescents, parents and therapists nested within a randomised controlled trial. Fine-tuning was declined in favour of frozen embeddings with weighted classifiers because positive instances were too sparse, an explicit acknowledgment that method must follow data size. Cross-validation was grouped by subject ID to prevent leakage across triplets, making this the only study in the section with documented subject-level leakage control. Some categories could not be trained at all.
	
	\subsection{Social media and platform corpora}
	
	Three studies construct corpora from public platform text, where the critical variable is not platform but label provenance.
	
	Qi et al. (2024)\cite{qi2024supervisedlearninglargelanguage} introduced SOS-HL-1K (1,249 posts, suicide risk) and SocialCD-3K (3,407 posts, cognitive distortion), crawled from the Zoufan blog on Weibo. The annotation protocol is the strongest in the section: three annotators with clinical psychology backgrounds worked double-blind, disagreements were resolved with a senior licensed psychologist, and agreement was quantified by both Fleiss' $\kappa$ and Krippendorff's $\alpha$. Supervised baselines were trained across five random seeds and LLMs evaluated over five runs at five temperature settings, with paired-samples $t$-tests at $\alpha$ = 0.05. The finding that supervised models still outperform LLMs on complex clinical classification carries weight because of that protocol.
	
	Nanda et al. (2024)\cite{10714257} used the Self-reported Mental Health Diagnoses corpus, where labels derive from users' own disclosures on Reddit rather than clinical assessment. Construction is a sequence of heuristics: regex cleaning and token pruning, condition keyword lists from DSM-5 headings extended with synonyms, misspellings and vernacular terms, cosine-similarity filtering at 0.5, and downsampling of controls to approximately 8,000 posts. Each step is individually defensible, but together they determine which users enter the corpus, and evaluation rests on a single English-language dataset without clinical validation.
	
	Sood et al. (2024)\cite{sood2024understandingstudentsentimentmental} created SMILE-College, 793 records of student narrative feedback from the College Pulse Student Voice Survey, annotated in two collaborative stages by a graduate student and a validator. The emergence of a "Mixed" sentiment category is treated as a substantive finding rather than a labelling inconvenience. No formal agreement statistic is reported.
	
	\needspace{8\baselineskip}
	\begin{center}
		\textbf{Table 2.1: Label provenance across the corpus-building studies}
		\vspace{2pt}
		
		\scriptsize
		\setlength{\tabcolsep}{3pt}\renewcommand{\arraystretch}{1.1}
		\begin{tabularx}{\textwidth}{|X|X|X|X|}
			\hline
			\textbf{Provenance tier} & \textbf{Mechanism} & \textbf{Studies} & \textbf{Agreement statistic reported} \\
			\hline
			Expert clinical annotation & Psychiatrist or psychologist coding of clinical or platform text & Xu et al. (2025a)\cite{Xu2025.01.03.24318117}; Qi et al. (2024)\cite{qi2024supervisedlearninglargelanguage}; Roy et al. (2026)\cite{10.1145/3805697} & Fleiss' $\kappa$ + Krippendorff's $\alpha$ (Qi); Cohen's $\kappa$ 0.74/0.72 (Roy); psychiatrist-reviewed test set (Xu) \\
			\hline
			Structured self-report instrument & PHQ-8 or BDI-II administered under protocol & Ohse et al. (2024)\cite{OHSE2024101663}; Sood et al. (2023)\cite{10385480}; Bucur (2024)\cite{bucur2024leveraging} & Not applicable \\
			\hline
			Trained non-clinical annotators & Graduate students, bilingual validators, relevance raters & Sood et al. (2024)\cite{sood2024understandingstudentsentimentmental}; Bucur (2024)\cite{bucur2024leveraging} & Majority-vote and unanimity gold standards (Bucur); none reported (Sood) \\
			\hline
			Platform self-disclosure & User-stated diagnosis matched by keyword patterns & Nanda et al. (2024)\cite{10714257} & None \\
			\hline
			Model-generated pseudo-label & LLM assigns training labels at scale & Aggarwal et al. (2024)\cite{10973522}; Hu et al. (2024)\cite{hu2024psycollmenhancingllmpsychological} & None (Aggarwal); human proofreading (Hu) \\
			\hline
		\end{tabularx}
	\end{center}
	
	\subsection{Synthetic and simulated corpora}
	
	Two studies generate rather than collect data, and both report fidelity problems that make them more valuable as cautionary evidence than as demonstrations.
	
	De Duro et al. (2025)\cite{DEDURO2025100170} produced CounseLLMe, 400 simulated counselling dialogues of 20 quips each, 200 English and 200 Italian, with a psychotherapist advising on prompt development. Evaluation reconstructs dialogues as Textual Forma Mentis Networks and profiles emotion using Plutchik lexicons, comparing against 24 human transcripts from HOPE with Kruskal-Wallis tests on network distances. Manual coding found role-swapping and dialogue interference in Italian-language generation, and the models failed to reproduce the anger and frustration characteristic of human patients.
	
	Bucur (2024)\cite{bucur2024leveraging} tested synthetic augmentation directly, generating 30 synthetic Reddit-style posts for each of 90 BDI-II responses (2,700 per model) using GPT-3.5 and Llama-3 8B, then testing whether these improved semantic retrieval from the eRisk 2023 corpus of roughly four million sentences. Relevance was established by top-$k$ pooling ($k$ = 50) with three annotators producing majority-voting and unanimity gold standards, evaluated by Average Precision, R-Precision, P@10 and NDCG@1000. The result is negative: generated queries were too specific and degraded retrieval relative to the original questionnaire responses. MPNet, fine-tuned for semantic search, outperformed MentalRoBERTa, inverting the expectation that domain pre-training dominates.
	
	\subsection{Annotation protocols and label provenance}
	
	Three studies treat the production of labels, rather than the collection of text, as the object of investigation.
	
	Roy et al. (2026)\cite{10.1145/3805697} built a diagnostic annotation layer over PRIMATE, prompting models to identify spans evidencing PHQ-9 and GAD-7 criteria with outputs constrained to structured formats. Ground truth was built in two stages: GPT-4o generated candidate annotations and three expert clinicians validated a subset, achieving Cohen's $\kappa$ of 0.74 for PHQ-9 and 0.72 for GAD-7. The datasets were released. Fine-tuning proved difficult, with the base MentaLLaMA model failing zero-shot tasks by repeating its input verbatim.
	
	Aggarwal et al. (2024)\cite{10973522} interrogate label provenance directly, comparing four annotation sources on identical Reddit data, SMHD dictionary matching against Llama3, MentaLLaMA and Samantha-Mistral, then fine-tuning downstream classifiers on each. Two findings matter. The LLM annotators exhibited severe severity inflation, with the majority of posts labelled Severe and one model producing no Mild labels. Classifiers trained on LLM pseudo-labels showed large train-validation gaps (95.8\% against 70.31\% for Llama3), which the authors read as overfitting to unreliable labels. Together this constitutes empirical evidence that LLM-generated training labels are not a neutral substitute for expert annotation.
	
	Hu et al. (2024)\cite{hu2024psycollmenhancingllmpsychological} constructed a large Chinese psychological training corpus, 155k single-turn QA pairs, 11k multi-turn dialogues and 10k textbook knowledge items, using a three-step pipeline of generation, evidence judgment and refinement. The evidence-judgment stage extracts matching source segments for each generated sentence, providing a grounding check against unsupported generation. Human proofreaders validated the output. The authors note potential generation bias from using LLMs to compile supervised fine-tuning data, the concern Aggarwal et al. demonstrate empirically.
	
	\subsection{Appraisal of the data foundations}
	
	Clinical data is the exception: four of twelve studies use data generated in a clinical encounter, and only Xu et al. exceeds a thousand participants, so claims about clinical performance rest on a narrow base. Agreement statistics are rarely reported, with three of twelve giving a formal inter-annotator statistic and the remaining nine using single annotators, unquantified multi-stage validation, or labels derived from self-report or model output. Where a corpus becomes a shared benchmark, unquantified annotation quality propagates silently into every study that uses it.
	
	Class imbalance is endemic and unevenly handled, named as limiting by four studies, with treatment ranging from SMOTE oversampling and weighted loss with subject-grouped folds to nothing at all. Leakage control is almost entirely undocumented, with only Xin et al. grouping cross-validation folds by subject. Given that several corpora contain multiple documents per participant, this is a material gap that inflates reported performance in ways no metric will reveal.
	
	\needspace{8\baselineskip}
	\begin{center}
		\textbf{Table 2.2: Corpus scale and validation characteristics}
		\vspace{2pt}
		
		\scriptsize
		\setlength{\tabcolsep}{3pt}\renewcommand{\arraystretch}{1.1}
		\begin{tabularx}{\textwidth}{|X|X|X|X|X|X|}
			\hline
			\textbf{Study} & \textbf{Corpus} & \textbf{Scale} & \textbf{Ground truth source} & \textbf{Imbalance handling} & \textbf{Leakage control} \\
			\hline
			Ohse et al. (2024)\cite{OHSE2024101663} & KID (+ DAIC/E-DAIC) & Semi-clinical interviews & PHQ-8 self-report & Not reported & Not reported \\
			\hline
			Sood et al. (2023)\cite{10385480} & I-DAIC & 626 documents & PHQ-8 self-report & Acknowledged, unresolved & Not reported \\
			\hline
			Xu et al. (2025a)\cite{Xu2025.01.03.24318117} & SMHC dialogues & 1,160 outpatients, $\sim$15,000 min & Psychiatrist EMR + 50-case expert test set & SMOTE + 10-fold CV & Not reported \\
			\hline
			Xin et al. (2024)\cite{xin2024outcomes} & IMPACT-ME & 34 subjects, 102 interviews & Qualitative content coding & Weighted loss & Subject-grouped 4-fold CV \\
			\hline
			Qi et al. (2024)\cite{qi2024supervisedlearninglargelanguage} & SOS-HL-1K; SocialCD-3K & 1,249 + 3,407 posts & 3 clinical annotators + adjudicator & Not reported & Not reported \\
			\hline
			Nanda et al. (2024)\cite{10714257} & SMHD subset & $\sim$8k control posts & Platform self-disclosure & Downsampling & Not reported \\
			\hline
			Sood et al. (2024)\cite{sood2024understandingstudentsentimentmental} & SMILE-College & 793 records & 2-stage student annotation & Not reported & Random 75/5/20 split \\
			\hline
			De Duro et al. (2025)\cite{DEDURO2025100170} & CounseLLMe & 400 dialogues & Synthetic (HOPE as reference) & Not applicable & Not applicable \\
			\hline
			Bucur (2024)\cite{bucur2024leveraging} & eRisk 2023 + synthetic & $\sim$4M sentences; 2,700 synthetic/model & 3 annotators, majority + unanimity & Not applicable & Not applicable \\
			\hline
			Roy et al. (2026)\cite{10.1145/3805697} & PRIMATE + released sets & Randomised subsets & GPT-4o proposal + 3 clinician validation & Not reported & Not reported \\
			\hline
			Aggarwal et al. (2024)\cite{10973522} & Reddit & 5,000 posts & LLM pseudo-labels (4 sources compared) & 5-fold CV & Not reported \\
			\hline
			Hu et al. (2024)\cite{hu2024psycollmenhancingllmpsychological} & Chinese corpus & 155k + 11k + 10k items & LLM generation + evidence judgment + proofreading & Not reported & Not reported \\
			\hline
		\end{tabularx}
	\end{center}
	
	Both studies that generated training or query data report fidelity failures, and the study that used model-generated labels at scale documented systematic severity inflation. Against this, Hu et al. show that an explicit evidence-judgment stage can impose grounding discipline on generated data. The implication carried into Sections 3 and 6 is that synthetic data in this domain appears to require a verification stage as a matter of course rather than as an enhancement.
	\section{Model Adaptation Paradigms}
	
	Fourteen studies (21.2\%), the largest grouping in this review, contribute a technique for adapting a general-purpose language model to a mental health task rather than a corpus, benchmark or clinical evaluation. Four strategies are represented, best understood as positions on a trade-off between the cost of changing model weights and the cost of controlling behaviour at inference. Prompt engineering changes nothing and controls at inference; fine-tuning changes weights and accepts the compute and data cost; retrieval augmentation leaves weights untouched but constrains generation against an external evidence store; orchestration composes several models or tools without adapting any. A fifth, smaller line abandons generation and treats the model as a feature extractor.
	
	\subsection{Prompt engineering and structured reasoning}
	
	Patil and Gedhu (2026)\cite{patil2026cognitivementalllm} conducted the cleanest ablation of reasoning strategies here, comparing Chain-of-Thought, Self-Consistency CoT, Few-Shot CoT and Tree-of-Thought on a single reasoning-oriented model (o3-mini) across five Reddit benchmarks. Holding the model constant while varying only the reasoning scaffold isolates the contribution of the strategy in a way multi-model comparisons cannot. The metric suite extends beyond accuracy and F1 to Matthews Correlation Coefficient, Quadratic Weighted Kappa, Mean Absolute Error and both ROC and PR AUC, appropriate choices for imbalanced ordinal tasks. Findings are differentiated: Few-Shot CoT performed best on multi-class tasks, plain CoT and SC-CoT on binary tasks, and Tree-of-Thought was inconsistent throughout. No inferential statistics are performed, so no claim of significant difference is licensed. Performance degrades in multi-class and long-text settings and remains sensitive to minor phrasing variation, itself a finding about the fragility of prompt-based adaptation.
	
	Kim et al. (2025)\cite{fpsyt.2025.1583739} developed LLM4CBT, a training-free prompt-alignment framework combining therapist persona, CBT technique examples and preferred clinical behaviours. Evaluation covers real therapist utterances from a merged HighQuality and HOPE corpus (366 dialogues, 7,669 utterances) against ground-truth act labels across 13 act types, plus simulated conversations in which elicited automatic thoughts are checked against ground truth embedded in generated patient profiles. Testing pacing against deliberately reluctant patients is a useful stress condition. The authors acknowledge a circularity risk, since the evaluation pipeline depends on automated LLM classifiers and therefore shares the failure modes of the system measured.
	
	Xu et al. (2024)\cite{Xu_2024} span this subsection and the next, and are placed here because the prompt taxonomy is the most systematic in the corpus. The zero-shot template decomposes into four components, with three enhancement strategies compared against a basic condition: context enhancement, mental health enhancement, and their combination. Few-shot and chain-of-thought variants extend the design, and crossing this axis with instruction fine-tuning makes this the only study to quantify the prompting-versus-tuning trade-off on identical tasks. Compact fine-tuned models (Mental-Alpaca 7B, Mental-FLAN-T5 11B) match task-specific Mental-RoBERTa and exceed GPT-4 in zero- and few-shot settings, which is the strongest evidence here that prompt-based adaptation of large proprietary models is not the efficient frontier for classification.
	
	\subsection{Supervised and parameter-efficient fine-tuning}
	
	Five studies update model weights, and their most striking common feature is the compute range they span, which functions as an implicit statement about who can participate in this research.
	
	Lai et al. (2023)\cite{lai2023psyllmscalingglobalmental} built Psy-LLM through domain pre-training on 2.85 GB of crawled Chinese psychology data over 100,000 iterations, then fine-tuning on 56,000 PsyQA pairs, on a single V100. Evaluation is dual: intrinsic metrics (perplexity, ROUGE-L, Distinct-1/2) alongside human evaluation in which six psychology-faculty students scored 200 pairs on helpfulness, fluency, relevance and logic. The authors concede that outputs fall short of human quality and that the unidirectional architecture constrains contextual memory.
	
	Maurya et al. (2025)\cite{Maurya2025} compared four lightweight sequence-to-sequence models fine-tuned on 20,500 QA pairs from four public counselling resources, with a 70/30 split on a single T4. Reporting hardware metrics alongside quality metrics makes deployability an explicit evaluation axis. The quality metrics are the weakness: ROUGE, BLEU-4 and perplexity measure $n$-gram overlap and are insensitive to empathy or clinical appropriateness, and the human assessment was conducted by a single non-expert evaluator, the thinnest in the corpus.
	
	George et al. (2024)\cite{10774623} fine-tuned Mistral-7B Instruct v0.1 with QLoRA on the counsel-chat dataset using a single RTX 3090. Evaluation departs from $n$-gram metrics in favour of PsychoBench, assessing Emotional Intelligence Scale and Empathy Scale scores against human reference averages from prior literature, with F-tests and $t$-tests. Applying instruments designed for humans to a model is of contested construct validity, but it is a more defensible proxy for counselling quality than BLEU.
	
	Gumilang et al. (2024)\cite{Gumilang2024DevelopmentOA} pursued cultural rather than architectural adaptation through a three-stage pipeline: global pre-training on Counsel Chat, expert-guided retraining on 1,800 locally sourced Indonesian entries, and L2 regularisation to counter overfitting on the small local corpus. Evaluation combines training and validation accuracy tracking with User Acceptance Testing among 58 students on Likert-scaled ease of use, responsiveness and satisfaction. Accuracy curves and satisfaction ratings jointly say little about clinical quality, and the system has no crisis detection or intervention protocol.
	
	Bookanakere et al. (2024)\cite{10743711} fine-tuned GPT-2 (1.5B) quantised to 4 bits with a documented configuration chosen to fit strict resource constraints. The pipeline transforms tabular psychiatric lifestyle data into narrative paragraphs before training, evaluated on a held-out partition of 1,000 cases. Two conceded limitations are material: quantisation entails precision loss, and the source dataset is synthetic, so tabular-to-text transformation may not resemble clinical language.
	
	\needspace{8\baselineskip}
	\begin{center}
		\textbf{Table 3.1: Compute footprint and evaluation instrument across fine-tuning studies}
		\vspace{2pt}
		
		\scriptsize
		\setlength{\tabcolsep}{3pt}\renewcommand{\arraystretch}{1.1}
		\begin{tabularx}{\textwidth}{|X|X|X|X|X|X|}
			\hline
			\textbf{Study} & \textbf{Base model} & \textbf{Method} & \textbf{Hardware} & \textbf{Primary evaluation} & \textbf{Human evaluation} \\
			\hline
			Lai et al. (2023)\cite{lai2023psyllmscalingglobalmental} & PanGu 350M / WenZhong 110M & Domain pre-training + SFT & V100 32GB & Perplexity, ROUGE-L, Distinct-1/2 & 6 psychology students, 200 pairs \\
			\hline
			Xu et al. (2024)\cite{Xu_2024} & Alpaca 7B / FLAN-T5 11B & Multi-task instruction FT & 8 $\times$ A100 80GB & Balanced accuracy, cross-dataset transfer & Not reported \\
			\hline
			George et al. (2024)\cite{10774623} & Mistral-7B Instruct & QLoRA & 1 $\times$ RTX 3090 24GB & PsychoBench EIS/Empathy, F-test, $t$-test & Human norms from literature \\
			\hline
			Maurya et al. (2025)\cite{Maurya2025} & T5, FLAN-T5, BART, GODEL (60--220M) & SFT & 1 $\times$ T4 16GB & ROUGE, BLEU-4, perplexity + hardware metrics & 1 non-expert evaluator \\
			\hline
			Gumilang et al. (2024)\cite{Gumilang2024DevelopmentOA} & GPT-4-based & SFT + L2 regularisation & Not reported & Train/validation accuracy & 58 students (UAT, Likert) \\
			\hline
			Bookanakere et al. (2024)\cite{10743711} & GPT-2 1.5B & 4-bit quantised FT & Edge-constrained & Accuracy, precision, recall, F1 & Not reported \\
			\hline
		\end{tabularx}
	\end{center}
	
	\subsection{Retrieval-augmented and knowledge-grounded architectures}
	
	Three studies leave weights untouched and constrain generation against external evidence. All three invoke hallucination reduction as motivation, and they differ sharply in whether they substantiate it.
	
	Kharitonova et al. (2025)\cite{Kharitonova15062025} built a retrieval system over the Clinical Practice Guidelines of the Spanish National Health System, using semantic search on paragraph embeddings with a negative-constraint prompt confining answers to retrieved content, isolating reasoning ability from parametric knowledge. Twenty expert-authored questions were rated by a specialist clinician on three Likert scales with independent verification by a second medical expert. The Veracity scale is the notable instrument: scored from $-1$ to $+1$ and explicitly penalising non-answers, it prevents a system scoring well by refusing to respond, a failure mode abstention-based designs invite. The claim that hallucination was eliminated by the closed-database design is an architectural argument rather than a measurement. Open-source LLaMA models answered more often and more accurately than GPT-3, which declined 25\% of ADHD questions.
	
	Guo and Guo (2024)\cite{10958051} integrated GraphRAG over a knowledge graph constructed from DSM-5 and ICD-11 criteria with EmoLLM for symptom extraction, requiring no additional training. Consultation dialogues are segmented into $k$-round groups ($k$ = 6), candidate symptoms extracted, then validated and deduplicated against the knowledge graph before a CoT-based diagnostic step. That validation stage is the most concrete hallucination-control mechanism in the section: rather than trusting retrieval to suppress fabrication implicitly, it re-checks each extracted symptom against a structured clinical ontology. The evaluation is appropriately ablated across zero-shot model comparison, backbone substitution, segmentation length and component removal, though it rests on a single small dataset (MMDA, 524 samples).
	
	Bhanu Sree et al. (2024)\cite{Sree2024RetrievalAugmentedGB} combined retrieval augmentation with LangChain agents executing specialised tools in parallel, routing queries by source type. The evaluation is almost entirely operational: time to first token, tokens per second and latency per minute, with qualitative feedback on conciseness, contextual accuracy, helpfulness and relevance. Latency outliers up to 495 seconds are reported. No hallucination measurement is performed despite grounding in PubMed being the stated rationale, and the study relies on general pre-trained models without domain-specific updates.
	
	\subsection{Orchestration and hybrid pipelines}
	
	Two studies compose components without adapting any of them, and both illustrate an evaluation problem worth naming.
	
	Kamoji et al. (2024)\cite{10675718} built a two-stage pipeline in which an ensemble classifier predicts mental health status from a 41,000-entry survey dataset with 25 features, after which Gemini generates personalised natural-language insight. The architecture is arguably prudent, since the diagnostic decision stays with an interpretable model. Two problems limit what can be concluded. Reported accuracy near 99\% on self-report data is, as the authors concede, suggestive of overfitting. The LLM component is never evaluated for insight quality, factual accuracy or safety, despite the reports being intended for psychiatrist review.
	
	Singh et al. (2024)\cite{10427865} documented MindGuide, a LangChain composition using LLMChain with ConversationBufferMemory over GPT-4 at temperature 0.5 through a Streamlit interface. The architectural documentation is clear and reproducible. The evaluation consists of a qualitative walkthrough of four conversation stages, with no metrics, dataset, comparison condition or formal human evaluation, which the authors state as a limitation. This matters beyond the individual paper because the system is positioned as an early identification assistant for anxiety, depression and suicidal thoughts with no crisis escalation protocol, guardrail or safety evaluation reported. It is included here as evidence about the state of the field's publication standards.
	
	\subsection{Embedding-based adaptation}
	
	Luna-Jiménez et al. (2024)\cite{10601118} abandon generation and use the LLM as a frozen feature extractor. Last-layer 4,096-dimensional d-vectors from Llama-2-7b-chat and MentaLLaMA-chat-7B were averaged temporally, ranked by F-ANOVA filter and passed to conventional classifiers. Evaluation used 80/10/10 splits on Counsel-Chat and 7Cups with weighted F1 alongside 95\% confidence intervals, one of the few studies here to quantify uncertainty in its headline metric. Roughly 75\% of dimensions proved redundant, permitting large reductions in downstream cost, though selected dimensions show low overlap across datasets, so the ranking appears dataset-specific rather than transferable.
	
	\subsection{Appraisal of the adaptation literature}
	
	Evaluation quality is inversely related to engineering ambition. The studies that build the most, whether full pipelines, agentic orchestration or deployed interfaces, evaluate the least: Singh et al. report no metrics, Kamoji et al. never evaluate the generative component of a two-stage system, and Bhanu Sree et al. measure latency but not accuracy or hallucination. The narrowest studies evaluate best, with Patil and Gedhu reporting ten metrics across five datasets, Luna-Jiménez et al. reporting confidence intervals, and Kharitonova et al. designing a scoring scale specifically to close an abstention loophole. A reader of this literature should treat architectural novelty and evidential strength as largely independent.
	
	Generation quality is measured with instruments known to be inadequate. Four studies generating therapeutic text rely wholly or partly on ROUGE, BLEU or perplexity, which measure overlap with a reference and are insensitive to empathy, clinical appropriateness or harm. The alternatives visible here, PsychoBench scales, act-label distribution matching against human therapists, and expert Likert rating with a non-answer penalty, are better aligned to the construct, and none has been adopted widely.
	
	Human evaluation, where present, is thin, ranging from six psychology students rating 200 pairs down to a single non-expert evaluator and none at all in five studies, and no study here reports inter-rater agreement for it. Hallucination control is asserted architecturally more often than measured: all three retrieval studies invoke hallucination reduction, but one implements an explicit verification stage, one argues from closed-database design without measuring, and one measures nothing. This sets up the distinction drawn in Section 6 between hallucination mitigation and hallucination evaluation.
	
	Compute access shapes the method distribution. Three studies frame parameter-efficient or quantised adaptation as enabling deployment in low-resource settings, which is precisely where clinician shortages are most acute, so the most deployable systems are trained on the smallest and least curated corpora, compounding the data-quality concerns of Section 2.
	
	\needspace{8\baselineskip}
	\begin{center}
		\textbf{Table 3.2: Evaluation adequacy across the adaptation studies}
		\vspace{2pt}
		
		\scriptsize
		\setlength{\tabcolsep}{3pt}\renewcommand{\arraystretch}{1.1}
		\begin{tabularx}{\textwidth}{|X|X|X|X|X|X|}
			\hline
			\textbf{Adaptation strategy} & \textbf{Studies} & \textbf{Quantitative task metrics} & \textbf{Generation-quality instrument} & \textbf{Human evaluation} & \textbf{Hallucination measured} \\
			\hline
			Prompting & Patil, Kim, Xu & All three & Act-label match (Kim) & LLM-based (Kim) & No \\
			\hline
			Fine-tuning & Lai, Maurya, George, Gumilang, Bookanakere & All five & ROUGE/BLEU (Lai, Maurya); PsychoBench (George) & 4 of 5, quality varies & No \\
			\hline
			Retrieval-augmented & Kharitonova, Guo, Bhanu Sree & Guo only & Expert Likert (Kharitonova) & Kharitonova (2 experts); Bhanu Sree (informal) & Verification stage (Guo) \\
			\hline
			Orchestration & Kamoji, Singh & Kamoji (classifier only) & None & None & No \\
			\hline
			Embeddings & Luna-Jim\'enez & Yes, with 95\% CIs & Not applicable & No & Not applicable \\
			\hline
		\end{tabularx}
	\end{center}
	
	\section{Benchmark Construction and Automated Evaluation}
	
	Six studies (9.1\%) contribute an evaluation artefact: a benchmark, a scoring protocol, or a framework for deciding when automated evaluation can be trusted. They are treated before human and clinical evaluation because the recurring question in that literature is whether expert raters are necessary, and that question is answerable only once the automated alternative and its documented failure points have been characterised. Two of the six address it directly by evaluating LLM judges against human experts on identical material. Three groupings emerge, distinguished by the source of ground truth: examination-grounded benchmarks inherit it from professional licensing instruments, expert-anchored benchmarks generate it from clinician panels, and judge-reliability studies treat the relationship between the two as the object of measurement. Methodological rigour in this corpus is concentrated here: five of six report inferential statistics, three report formal reliability coefficients, and two quantify uncertainty in their headline estimates.
	
	\subsection{Examination-grounded competency benchmarks}
	
	Nguyen et al. (2025)\cite{nguyen-etal-2025-large} constructed CounselingBench from professional licensing standards: 1,612 multiple-choice questions across 138 clinical case studies with expert-written rationales. Twenty-two models were crossed with five core competencies under zero-shot, few-shot, chain-of-thought and self-consistency decoding. The critical design decision pairs each medical-specialised model with its own general-purpose counterpart, isolating the effect of domain fine-tuning from architecture, so the finding that medical models underperform their generalist bases by 4.2 percentage points is attributable rather than merely correlational. The mechanism is traced through manual annotation of reasoning errors on 100 responses by three expert annotators. Competency mapping used two psychiatrists with a licensed counsellor as tiebreaker, and paired $t$-tests and chi-square tests support the accuracy and error-distribution claims. Multiple-choice cannot capture empathy or therapeutic alliance, and a US licensing exam has limited applicability to the Global South.
	
	Xu et al. (2025b)\cite{Xu2025} built the first integrated Chinese benchmark, combining the Chinese Counselor Qualification Exam (744 single-choice and 200 multiple-choice items, 2023--2024) with translated social media diagnostic screening on Dreaddit (1,151 posts) and SDCNL (1,517 posts). The two-task structure separates theoretical knowledge from applied diagnostic performance and finds these dissociate, since models strong on examination content are not correspondingly strong on classification. Fifteen models spanning 1.5B to 671B parameters were queried through a programmatic API harness at default parameters, a choice that maximises reproducibility while not exploring decoding sensitivity. Class balance was maintained in SDCNL test sets, and models over-predict the anxiety class under uncertainty. Analysis is descriptive, with accuracy compared across parameter scales and release timelines and no inferential testing. Ground truth quality is uneven across sources, with professional examiners for the exam but crowdworkers for the two social media corpora.
	
	Hanss et al. (2025)\cite{Hanss2025} took 150 single-answer multiple-choice questions from a psychiatry review manual and administered each ten times to GPT-3.5, GPT-4 and GPT-4o, allowing up to fifteen attempts to obtain ten valid responses. The repetition is the point. Rather than reporting a single accuracy figure, the study computes response consistency via frequency variance across trials and tests it as a predictor of correctness, alongside separately elicited pre-answer self-reported confidence. Consistency predicts accuracy while self-reported confidence does not, which supplies a deployable proxy for reliability requiring no ground truth at inference time. Statistical treatment is the most careful in this subsection, with chi-square tests under Bonferroni correction at $\alpha$ = .01, $t$-tests for consistency and point-biserial correlations for predictors, and documented decoding parameters. Two threats are flagged candidly: the format cannot assess interactive dialogue, and the questions may already sit in model training data, a contamination risk intrinsic to benchmarks built from published materials.
	
	\needspace{8\baselineskip}
	\begin{center}
		\textbf{Table 4.1: Examination-grounded benchmark parameters}
		\vspace{2pt}
		
		\scriptsize
		\setlength{\tabcolsep}{3pt}\renewcommand{\arraystretch}{1.1}
		\begin{tabularx}{\textwidth}{|X|X|X|X|X|X|}
			\hline
			\textbf{Study} & \textbf{Instrument} & \textbf{Scale} & \textbf{Models} & \textbf{Repetition} & \textbf{Inferential statistics} \\
			\hline
			Nguyen et al. (2025)\cite{nguyen-etal-2025-large} & Professional licensing standards & 1,612 MCQs, 138 case studies & 22 (13 medical + matched generalists) & 4 prompting/decoding conditions & Paired $t$-tests; chi-square \\
			\hline
			Xu et al. (2025b)\cite{Xu2025} & Chinese Counselor Qualification Exam + 2 social corpora & 944 exam items; 2,668 posts & 15 (1.5B--671B) & Single pass, default parameters & None (descriptive) \\
			\hline
			Hanss et al. (2025)\cite{Hanss2025} & Psychiatry review manual & 150 MCQs $\times$ 10 trials & 3 (GPT family) & 10 trials per item & Chi-square (Bonferroni, $\alpha$=.01); $t$-tests; point-biserial $r$ \\
			\hline
		\end{tabularx}
	\end{center}
	
	\subsection{Expert-anchored benchmarks}
	
	Two studies generate ground truth from clinicians rather than inheriting it, converting that expense into evaluative depth unavailable to multiple-choice designs.
	
	Li et al. (2026)\cite{li2026counselbenchlargescaleexpertevaluation} assembled the largest clinical expert panel in the corpus. COUNSELBENCH-EVAL comprises 100 real patient questions from CounselChat with four responses each, three LLM-generated and one from a verified human therapist, rated blind by 100 licensed or professionally trained practitioners. Rating covers six dimensions with scales matched to the construct, and raters extracted supporting text evidence and wrote rationales, so judgements trace to specific spans rather than holistic impressions. Inter-rater reliability is reported as Krippendorff's $\alpha$, appropriate for the mixed scale types, and inference uses Wilcoxon signed-rank tests with Kruskal-Wallis and chi-square examining whether rater experience or time-on-task influenced ratings, a form of rater-effect control almost absent elsewhere. The second component is adversarial: ten clinicians authored 120 questions targeting six failure modes identified in the first phase, and five different practitioners evaluated 1,080 responses, a separation preventing self-marking. Hallucination is operationalised as Factual Consistency scoring validated against expert annotation, and safety as Medical Advice and Toxicity flags. A pilot excluded three domain-tuned models for high invalid-response rates, which is itself a finding. Limitations are single-turn interaction and reliance on scarce non-deidentified data.
	
	McBain et al. (2025a)\cite{MCBAIN2025} repurposed the Suicide Intervention Response Inventory (SIRI-2), 24 hypothetical patient remarks paired with two clinician responses each, as an LLM benchmark. The decision giving the study its force is prompting deliberately without engineering, using original SIRI-2 instructions only and zero-shot, to isolate baseline competency rather than optimised performance. Because SIRI-2 carries published norms, model scores are positioned directly on a human scale spanning expert suicidologists, clinical PhDs, master's-level counsellors and untrained K-12 school staff. Reporting the resulting placements is more interpretable than an accuracy percentage. Three independent research accounts prompted each model, supporting test-retest correlation. Analysis accounts for nested item structure through linear regression with z-score outlier detection at $|z| > 1.96$, conducted under STROBE reporting. The stated limitation is precise: the study measures evaluative competency, the model as referee, not active engagement with a person in crisis.
	
	\subsection{Judge reliability and the limits of automated evaluation}
	
	Badawi et al. (2026)\cite{badawi-etal-2026-trust} address when an LLM judge can substitute for a human expert. MentalBench-100k comprises 10,000 conversations from three sources and MentalAlign-70k comprises 70,000 ratings. From 1,000 curated contexts, ten responses each were generated, one human and nine model, and rated by four LLM judges and three human experts with graduate-level or licensed psychiatric backgrounds across seven attributes grouped into a Cognitive Support Score and an Affective Resonance Score.
	
	The statistical apparatus is the most developed in the corpus. Agreement is decomposed into consistency ICC(C,1) and absolute agreement ICC(A,1), a distinction that matters because a judge can rank responses correctly while being systematically miscalibrated, and only the second coefficient detects it. Systematic bias is reported directly as the LLM-human difference. Precision is quantified through 95\% bootstrap confidence interval widths from 1,000 nonparametric iterations, with two-way mixed-effects ANOVA variance decomposition and paired $t$-tests supporting inference across 28 judge-attribute combinations. Ratings were aggregated to model level to filter conversation-level noise, and traditional error metrics serve as baselines. The resulting Affective-Cognitive Agreement Framework delivers a conditional rather than blanket verdict, classifying where automated judgement is reliable and where human oversight remains mandatory. That conditional structure converts an ideological dispute about automated evaluation into a calibration problem. Limitations are English-only, single-turn dialogue, the computational cost of running multiple large judges, and expert rater availability, with three raters supporting 70,000 ratings.
	
	\subsection{Appraisal of the benchmark literature}
	
	Statistical rigour is markedly higher here than elsewhere, with five of six studies reporting inferential statistics against roughly one in five across the adaptation literature, three reporting formal reliability coefficients and two quantifying uncertainty explicitly. Xu et al. (2025b) is the sole exception, reporting descriptive accuracy only despite the largest model sweep in the subsection.
	
	Stochasticity is taken seriously, but unevenly: Hanss et al. administer ten trials per item and convert the variance into a reliability measure, McBain et al. use three independent instances, Li et al. fix and document decoding parameters, and Xu et al. run a single pass at default settings. Since identical prompts yield different answers across runs, single-pass results carry unreported variance and the field has not converged on a repetition standard.
	
	The multiple-choice format is the binding constraint on examination-grounded work, since performance cannot capture empathy, therapeutic alliance or interactive competence. Expert-anchored designs recover open-ended assessment at a cost visible in their sample sizes, 100 questions and 48 items respectively, and the trade-off between construct validity and scale remains unresolved. Contamination is acknowledged once and otherwise unaddressed, though the risk applies to every benchmark built from published instruments.
	
	Domain fine-tuning does not help, on two independent measurements. Nguyen et al. find medical-specialised models underperform matched generalist bases by 4.2 percentage points, and Li et al. excluded three domain-tuned models for high invalid-response rates. Combined with the finding in Section 3 that compact instruction-tuned models beat far larger prompted ones, the emerging picture suggests that task-directed adaptation succeeds where domain-directed adaptation does not.
	
	\needspace{8\baselineskip}
	\begin{center}
		\textbf{Table 4.2: Evaluation rigour across the benchmark studies}
		\vspace{2pt}
		
		\scriptsize
		\setlength{\tabcolsep}{3pt}\renewcommand{\arraystretch}{1.1}
		\begin{tabularx}{\textwidth}{|X|X|X|X|X|X|}
			\hline
			\textbf{Study} & \textbf{Ground truth source} & \textbf{Reliability statistic} & \textbf{Uncertainty quantified} & \textbf{Repetition} & \textbf{Rater-effect control} \\
			\hline
			Nguyen et al. (2025)\cite{nguyen-etal-2025-large} & Licensing exam key + 2 psychiatrists & Tiebreaker adjudication & No & Prompting conditions & Not reported \\
			\hline
			Xu et al. (2025b)\cite{Xu2025} & Exam key + crowdworker labels & None & No & Single pass & Not applicable \\
			\hline
			Hanss et al. (2025)\cite{Hanss2025} & Published review manual key & Response consistency (variance) & Bonferroni $\alpha$ = .01 & 10 trials/item & Not applicable \\
			\hline
			Li et al. (2026)\cite{li2026counselbenchlargescaleexpertevaluation} & 100 licensed practitioners, blinded & Krippendorff's $\alpha$ & No & Fixed decoding & Kruskal-Wallis/chi-square on experience and time \\
			\hline
			McBain et al. (2025a)\cite{MCBAIN2025} & Expert suicidologist norms + historical cohorts & Test-retest correlation & z-score outlier detection & 3 instances & Nested-structure regression \\
			\hline
			Badawi et al. (2026)\cite{badawi-etal-2026-trust} & 3 psychiatric experts, 70k ratings & ICC(C,1) and ICC(A,1) & 95\% bootstrap CIs (1,000 iters) & 4 judges $\times$ 7 attributes & Two-way mixed-effects ANOVA \\
			\hline
		\end{tabularx}
	\end{center}
	
	Two studies here turn the field's evaluation apparatus on itself, testing whether LLM judges align with human experts on the same material. Rather than adopting LLM-as-judge for convenience and defending it informally, both quantify the conditions under which it holds. The distinction Badawi et al. draw between consistency and absolute agreement carries into Section 7, where it applies equally to the human rating protocols examined in Section 5.
	\section{Human-Centred and Clinical Evaluation Designs}
	
	Expert time is the scarcest resource in this literature, and most design decisions examined below are shaped by it. Fourteen studies (21.2\%) place human judgement or human experience at the centre of evaluation. They are ordered here along a validity ladder. Vignette studies maximise experimental control at the cost of ecological realism; expert-rater studies apply clinical judgement to real model outputs; controlled trials and prospective deployments observe systems in use; qualitative work studies naturalistic interaction without intervening; and population survey evidence describes what is happening at scale. Internal validity is highest at the top of this ladder and ecological validity at the bottom, and no study achieves both.
	
	\subsection{Vignette and standardised-case designs}
	
	Four studies present models with controlled clinical stimuli, multiplying small vignette sets through repetition to achieve analysable sample sizes.
	
	Levkovich (2025)\cite{ejihpe15010009} administered six published clinical vignettes to four LLMs twenty times each in parallel male and female versions, giving 960 evaluations. The comparison standard is exceptional: norms from 1,536 health professionals and 6,016 members of the public, inherited from the source studies. Open-ended diagnoses were scored by automated string-search filters, with chi-square and Fisher's exact tests and Cram\'er's V effect sizes. The gender-balanced design doubles as a bias probe. The authors are direct about the ceiling on their claims: hypothetical vignettes, no multi-turn dynamics, and no clinical validation.
	
	Schnepper et al. (2025)\cite{SCHNEPPER2025} built the tightest experimental design in the section, a randomised 2 $\times$ 2 crossing gender with sexual orientation across 30 published eating-disorder vignettes, yielding 120 variants evaluated in three rounds by each of two models, for 720 evaluations. Outcome instruments are validated psychometrics rather than ad-hoc scales, the RAND-36 Mental Composite Summary and the EDE-Q global score, rendering any bias effect interpretable in clinically familiar units. Analysis uses linear multilevel models with crossed random effects and intraclass correlation coefficients testing replicability, with temperature set to 0. Three limitations matter: the source vignettes were heavily gender-imbalanced, the custom temperature controls are unverified, and MentaLLaMA produced data of such low quality that the general-versus-specialised comparison is substantially weakened.
	
	Thotapalli et al. (2025)\cite{pcn5.70159} evaluated three GPT variants across 22 youth psychiatric emergency vignettes over four iterations (264 responses) against two advanced practice nurse practitioners, with three board-certified psychiatrists grading response quality. The reliability treatment is the strongest among the vignette studies, using quadratically weighted Cohen's $\kappa$ for clinician-LLM agreement, appropriate for ordinal triage categories where the distance between disagreements matters, alongside Fleiss's $\kappa$ for inter-model consistency. The substantive finding is an error asymmetry, zero false negatives against up to four false positive admissions, which is the safety-preferable direction but carries resource implications. Vignettes were drafted using ChatGPT and refined by the research team, a circularity the study should be read against.
	
	Lauderdale et al. (2025)\cite{Lauderdale2025EffectivenessOG} ran three models across four standardised veteran vignettes over ten trials each, scored against the Veterans Health Administration Risk Stratification Table and compared with 42 licensed mental health care providers. The instrument choice is the contribution: an operational structured professional judgement framework already in clinical use, with acute risk, chronic risk and treatment disposition rated separately. Separating acute from chronic risk is a distinction most studies collapse. Analysis used independent-groups $t$-tests, one-way ANOVA with Tukey's HSD and hierarchical regression. The authors concede that four vignettes cannot support counterbalancing of demographic factors, and note that existing VHA machine-learning risk models were validated largely on White male veterans.
	
	\needspace{8\baselineskip}
	\begin{center}
		\textbf{Table 5.1: Vignette study designs}
		\vspace{2pt}
		
		\scriptsize
		\setlength{\tabcolsep}{3pt}\renewcommand{\arraystretch}{1.1}
		\begin{tabularx}{\textwidth}{|X|X|X|X|X|X|}
			\hline
			\textbf{Study} & \textbf{Vignettes} & \textbf{Repetition} & \textbf{Total observations} & \textbf{Human comparator} & \textbf{Reliability statistic} \\
			\hline
			Levkovich (2025)\cite{ejihpe15010009} & 6 ($\times$ male/female) & 20 per version & 960 & 1,536 professionals + 6,016 public (published norms) & Cram\'er's V effect sizes \\
			\hline
			Schnepper et al. (2025)\cite{SCHNEPPER2025} & 30 $\rightarrow$ 120 variants (2$\times$2) & 3 rounds & 720 & None (published literature source) & ICC for replicability \\
			\hline
			Thotapalli et al. (2025)\cite{pcn5.70159} & 22 & 4 iterations & 264 & 2 nurse practitioners + 3 psychiatrists & Weighted Cohen's $\kappa$; Fleiss's $\kappa$ \\
			\hline
			Lauderdale et al. (2025)\cite{Lauderdale2025EffectivenessOG} & 4 & 10 trials & 120 & 42 licensed providers & Not reported \\
			\hline
		\end{tabularx}
	\end{center}
	
	\subsection{Expert-rater comparative studies}
	
	Settanni et al. (2025)\cite{SETTANNI2025116583} evaluated 425 Italian-language Reddit posts, filtered from 4,462, across three models under two prompt conditions, producing 2,550 ratings against a gold standard from two licensed clinical psychologists. Models output only a numerical urgency score on an adapted Mental Health Triage Scale customised for informal social media register. The evaluation reports continuous agreement alongside binary classification metrics after dichotomising at the urgent threshold, with repeated stratified 5-fold cross-validation over ten repetitions and temperature fixed at 0. Weighted Cohen's $\kappa$ quantifies clinician inter-rater reliability. As in Thotapalli et al., the dominant error is over-triage, characterised explicitly as the safety-relevant failure mode.
	
	Cui et al. (2025)\cite{fpsyt.2025.1634714} developed a suicide intervention chatbot around a five-stage crisis protocol grounded in intervention manuals and the ACT model, deployed via a web platform with documented decoding parameters. Twenty psychology professionals conducted guided therapy sessions and completed an eight-statement questionnaire across six dimensions on 7-point Likert scales, including a dedicated Safety and Privacy dimension. Analysis is descriptive, with no comparison condition and no inter-rater statistic, which limits the design to establishing expert acceptability rather than efficacy.
	
	Eberhardt et al. (2025)\cite{Eberhardt2025} is, by statistical standards, the most complete study in the corpus. From 1,131 video-recorded therapy sessions across 155 outpatients it develops an automated rating scale for latent patient engagement through a full psychometric pipeline, with 120 candidate items pre-selected on distributional properties, factor loadings and item-total correlations, then validated against independent process and outcome scales. Reliability is Cronbach's $\alpha$ = 0.952, McDonald's $\omega$ = 0.953 and average variance extracted = 0.715, with multilevel confirmatory factor analysis giving CFI = 0.968, TLI = 0.956, SRMR = 0.022 and RMSEA = 0.108, the last exceeding conventional thresholds and reported rather than suppressed. Overfitting is guarded against by 1,000-resample bootstrap optimism correction and 3-fold cross-validation, and transcription quality is quantified, a step almost no other study using transcribed speech takes. The model runs locally on secure institutional servers so patient data never leaves the institution. The trade-off is cost and a patient-level sample of 155 that is modest for scale validation.
	
	\subsection{Controlled trials and prospective deployment}
	
	Campellone et al. (2025)\cite{info:doi/10.2196/67365} provides the corpus's only randomised controlled trial: 160 adults randomised 1:1 to a generative AI arm ($n$ = 81) or an identical rules-based comparator ($n$ = 79) over two weeks, double-blind and decentralised. The contribution is the safety architecture and its verification. Guardrails operate at four tiers: a classifier that blocks concerning language and routes to helplines before any model call; an embedding-based off-topic classifier; prompt-level clinical constraints; and output-level structure and banned-word checks. Readiness was tested pre-trial against 42 simulated patient personas, and all 2,207 generated completions were manually annotated post-trial by two coders with principal-investigator resolution, verifying guardrail success exhaustively rather than by sampling. Analysis is descriptive, since the study was explicitly not powered for between-group inference and did not adjust for baseline severity.
	
	Patias et al. (2024)\cite{10845582} describe the me\_HeLi-D project, integrating GPT-4 into digital modules for adolescents aged 12--15 covering early screening from journal entries, CBT and mindfulness interventions, interactive lessons and conversational chatbots. Planned evaluation is mixed-methods, combining pre/post surveys using the Mental Health Literacy Scale and Strengths and Difficulties Questionnaire with interaction logs, focus groups and behavioural tracking. Data protection is specified concretely. The study is ongoing and its conclusions rest on preliminary data, so it contributes a protocol rather than findings. It is included because it is the only prospective work with minors in the corpus, and because automated screening of adolescent journal entries raises consent questions the paper acknowledges but does not resolve.
	
	\subsection{Qualitative, ethnographic and naturalistic studies}
	
	Four studies observe interaction rather than staging it, and collectively they supply the only evidence about what these systems do outside a research protocol.
	
	Wang et al. (2025b)\cite{wang2025evaluatingllmpoweredchatbotcognitive} structured evaluation in three stages with deliberate separation of roles: co-design with five mental health professionals averaging fifteen years of experience, a field user study with 19 participants producing 95 transcripts, and transcript review by four different professionals who had not participated in the design. Reviewers used a think-aloud protocol over randomised dialogue logs, analysed thematically with iterative consensus-building. Separating designers from evaluators prevents self-marking, a control absent from most system papers in Section 3. Thematic codes name specific clinical failure modes including toxic positivity, leading questions and insufficient Socratic depth.
	
	Song et al. (2025)\cite{song2025typingcureexperienceslarge} conducted semi-structured interviews with 21 participants recruited purposively and by snowball sampling from specialised online communities, with deliberate diversity across Indian, Nigerian, Brazilian and Central Asian participants. This is the most globally distributed sample in the corpus, and the choice is methodologically motivated, since general-purpose models carry embedded linguistic and cultural bias and a Western-only sample would systematically miss cultural mismatch. The analysis formalises therapeutic alignment as a multidimensional construct for assessing whether human-AI interaction supports long-term psychological healing. Limitations are a small non-clinical sample, retrospective self-report and no clinical outcome measurement.
	
	Iftikhar et al. (2025)\cite{Iftikhar2025HowLC} ran an 18-month ethnography combining 110 self-counselling sessions conducted by seven trained peer counsellors, who met in 60 weekly focus groups to evaluate LLM behaviour, with 27 simulated multi-turn sessions independently and blindly reviewed by three licensed clinical psychologists. The longitudinal design is unique in the corpus. Inductive coding refined 41 initial codes to 15 ethical risks across five themes, each mapped to formal professional standards. Mapping to codified professional ethics rather than author intuition gives the resulting framework normative force. No quantitative analysis is reported, and the study covers only prompted rather than fine-tuned models on a single CBT-focused platform, with all psychologists from one country.
	
	Moore et al. (2026)\cite{moore2026characterizingdelusionalspiralshumanllm} analysed real chat logs from 19 users reporting psychological harm, 391,562 messages across 4,761 conversations. The pipeline is technically careful: transcripts de-identified with Presidio and Faker, tree-structured chat histories linearised into ancestral chains, and 28 codes across five categories applied by an LLM annotator with confidence thresholds raised for high-risk codes. Validation is layered, with a random 560-message sample annotated by seven co-authors including a board-certified psychiatrist, plus full manual verification of every suicidality and violence code. Agreement is reported as Cohen's $\kappa$ for LLM-human and Fleiss' $\kappa$ for inter-annotator agreement. Conversational dynamics are modelled through per-message regressions on log-transformed remaining conversation length with participant-clustered standard errors. The study is candid that agreement is moderate and prone to false positives on subjective codes, the sample is small and self-selected, and no objective clinical outcomes exist.
	
	\subsection{Population survey evidence}
	
	Stade et al. (2026)\cite{Stade2026} surveyed 1,871 US adults between August and October 2024 using stratified sampling across age, sex and race/ethnicity to approximate national demographics, with hypotheses preregistered. All respondents answered demographic, psychopathology and treatment-use items; self-identified users additionally answered questions on how, why and for what they use LLMs, plus perceptions of usefulness, therapeutic alliance and the importance of memory. The headline estimate, 24\% of respondents using LLMs for mental health and extrapolating conservatively to 14--18 million US adults, is the only population-scale figure in the corpus, and the linkage to treatment access barriers grounds the literature in a demand context model-evaluation studies cannot supply. Limitations are inherent to the design: cross-sectional self-report supports no causal inference, coverage is single-country and single-window, and screening measures substitute for diagnoses. This study's extraction is derived from optical character recognition, since the source document carries no text layer, and should be verified against the original before citation.
	
	\subsection{Appraisal of the human evaluation literature}
	
	Reliability reporting is stronger here than anywhere else in the corpus, with six of fourteen studies giving a formal agreement or reliability coefficient against three of twelve in Section 2 and none in Section 3, the clearest evidence that measurement discipline tracks proximity to clinical practice. Expert availability sets sample size rather than the research question: two clinicians establish a gold standard for 2,550 ratings, three psychologists review 27 sessions, and four reviewers cover 95 transcripts. Vignette studies work around this by inheriting published norms, which is efficient but ties findings to whatever the source studies measured. Repetition is standard here, with all four vignette studies administering multiple trials and two fixing temperature at 0, so the stochasticity problem left unresolved in Section 4 is routinely handled.
	
	The prospective evidence base is one underpowered trial: two weeks, $n$ = 160, explicitly not powered for between-group inference, run by the product developer, and excluding participants with recent suicidality, precisely the population whose safety is most at issue. No study in this corpus reports clinical outcomes from a powered trial, so every efficacy claim rests on proxy measures. Conservative bias appears independently in two triage studies, both finding models erring toward escalation, which is the safety-preferable direction and contrasts with Section 6, where studies probing crisis handling find models escalating too slowly or not at all.
	
	\needspace{8\baselineskip}
	\begin{center}
		\textbf{Table 5.2: Position on the validity ladder}
		\vspace{2pt}
		
		\footnotesize
		\setlength{\tabcolsep}{3pt}\renewcommand{\arraystretch}{1.1}
		\begin{tabularx}{\textwidth}{|X|X|X|X|X|}
			\hline
			\textbf{Design} & \textbf{Studies} & \textbf{Internal validity} & \textbf{Ecological validity} & \textbf{Reliability statistic reported} \\
			\hline
			Vignette / standardised case & Levkovich, Schnepper, Thotapalli, Lauderdale & High (controlled stimuli, repetition) & Low (hypothetical, single-turn) & 3 of 4 \\
			\hline
			Expert rating of real outputs & Settanni, Cui, Eberhardt & Moderate--high & Moderate & 2 of 3 \\
			\hline
			Controlled trial / deployment & Campellone, Patias & Moderate (randomised but underpowered) & High & 0 of 2 \\
			\hline
			Qualitative / naturalistic & Wang, Song, Iftikhar, Moore & Low (no control condition) & Highest (real use, real harm) & 1 of 4 \\
			\hline
			Population survey & Stade & Not applicable (descriptive) & High (national sample) & Not applicable \\
			\hline
		\end{tabularx}
	\end{center}
	
	The studies with the strongest ecological validity are precisely those that can make no causal claims, while the study with the most rigorous psychometrics measures a single construct on a single sample. The field has not yet produced a design combining clinical realism with inferential power.
	
	\section{Risk, Safety and Assurance Evaluation}
	
	Assurance evaluates everything preceding it. A corpus can encode bias, an adaptation method can introduce fabrication, a benchmark can miss the failure that matters, and a clinical evaluation can be underpowered to detect harm. Ten studies (15.2\%) treat those failures as the object of measurement. The section has a dual structure: ten studies have risk evaluation as their primary design and are treated individually, but risk fields are populated far more widely, with hallucination methods appearing in 37 of 66 studies, bias in 58 and safety in 61. The corpus-wide extraction supports a claim the individual studies cannot, that high nominal coverage conceals a much smaller body of actual measurement.
	
	\subsection{Hallucination and factual grounding}
	
	Two studies operationalise hallucination as a measurable quantity, and they define it in incompatible ways.
	
	Linardon et al. (2025)\cite{LINARDON2025} converts hallucination into a verifiable outcome, citation fabrication. Six literature reviews of roughly 2,000 words were generated in a 3 $\times$ 2 factorial design crossing disorder familiarity with prompt specificity. All 176 citations were manually verified across Google Scholar, Scopus, PubMed and WorldCat, then classified as fabricated, real-with-errors or accurate, with error subtypes recorded down to author, year, title, journal, volume, issue, page and DOI. Two-tailed chi-square tests of independence were used with pairwise comparisons. The factorial structure gives the design its power: 19.9\% of citations were fabricated overall, but fabrication concentrated in the less-visible disorders (28\% and 29\% against 6\% for major depressive disorder). Hallucination therefore appears conditional on informational terrain rather than a fixed model property, a result no single-condition study could produce. Limitations are narrow scope: one model, one output per condition, and no test of mitigation.
	
	Kim et al. (2026)\cite{kim-etal-2026-kind} define a construct the factual literature does not capture. Affective hallucination denotes emotionally immersive responses evoking false social presence, a model simulating emotional capacity it does not possess. Prompts from five mental-health subreddits produced AHaBench (500 prompts) and AHaPairs (5,000 preference pairs), assessed on a 7-point rubric across Emotional Enmeshment, Illusion of Presence and Fostering Overdependence. The automated judge is validated against clinical ground truth, with two licensed psychiatrists rating outputs and human-GPT-4o agreement at $r$ = 0.85 against inter-human agreement of $r$ = 0.95. The study supplies mitigation as well as diagnosis, comparing DPO alignment against SFT, SFT+DPO and few-shot baselines across model sizes from 7B to 72B, and verifying on MMLU, GSM8K and ARC that alignment does not degrade general reasoning. Testing for capability regression after a safety intervention is rare here and arguably should be standard. Limitations are Reddit-sourced data with demographic skew, single-turn evaluation, and potential author annotation bias.
	
	\subsection{Bias and stigma auditing}
	
	Yeo et al. (2025)\cite{Yeo2025SociodemographicBiasMentalHealthAI} built a closed-loop agent-to-agent design of unusual care. A GPT-4 conversational agent interacted with GPT-3.5 Digital Standardised Patients across 97 demographic combinations of age, sex, race and income, each tested four to five times alongside controls, producing 449 transcripts and 4,502 agent responses. The critical control is asymmetric information: simulated patients remained agnostic to their own sociodemographic traits while the agent was informed of them through a separate channel, so any linguistic difference must originate with the agent. Measurement used LIWC-22 across 13 primary and 30 exploratory markers, deliberately avoiding human raters to eliminate researcher subjectivity. The statistical treatment is the most thorough of any bias study here, spanning normality testing, subgroup comparison, multivariate regression, Joinpoint regression for piecewise trend analysis of tone across ten sequential dialogue bins, and group-based trajectory modelling. Analysing tone trajectory within conversations, rather than aggregate difference, captures bias that emerges over an interaction. Limitations are simulated patients, a single English-language model, and LIWC's insensitivity to subtler implicit bias.
	
	Moore et al. (2025)\cite{Moore_2025} examined stigma and clinical appropriateness together. A mapping review of ten clinical guideline documents yielded 17 features of good therapy, used both to construct a steel-man system prompt, giving models their best chance rather than a strawman, and as the evaluation rubric. The first experiment applied modified National Stigma Studies vignettes, 72 vignettes across four conditions with 14 questions, measuring willingness to interact socially with the described character. The second tested responses to ten acute stimuli including suicidal ideation and delusions, with and without transcript context. Statistical rigour is high, with bootstrapped 95\% confidence intervals and Bonferroni-corrected $t$-tests and z-tests of proportion, and inter-rater reliability reaches Fleiss' $\kappa$ = 0.96, the highest in the corpus. Sixteen human therapists supplied baseline responses. The safety criterion is concrete: failing to challenge delusions, or supplying suicide-enabling information, is classified as unsafe.
	
	\subsection{Crisis-safety and escalation testing}
	
	McBain et al. (2025b)\cite{McBain2025LLMSuicideRisk} conducted the largest safety audit by volume. Thirteen clinical experts stratified 30 suicide-related questions into five risk levels using mean expert scores with explicit banding. Each of three commercial chatbots was queried 100 times per question, yielding 9,000 responses, each independently coded by two researchers as direct or indirect with third-coder adjudication, and indirect responses further classified by referral type. Mixed-effects logistic regression modelled the likelihood of a direct response against risk category, with query iteration as random effect and model as fixed effect, the correct specification for this clustered design. Stratifying refusal behaviour across five clinician-calibrated levels rather than a binary safe/unsafe split exposes inconsistency at intermediate risk that a binary framing would hide.
	
	Heston (2023)\cite{PMID:38111813} is the smallest study here and produces one of the most consequential findings. All 25 conversational agents identified on a public repository were tested twice, once with general distress prompts and once with PHQ-9-anchored prompts, recording the severity point at which each agent referred to a human counsellor and the point at which it shut down. Anchoring escalation to PHQ-9 makes slow escalation quantifiable in clinically interpretable units rather than impressionistic. The distinctive move is continuing past shutdown to test restartability, revealing that 22 of 25 safety-triggered agents resume normal conversation once distress severity drops, a loophole a single-pass test cannot detect and evidence that safety shutdowns are not durable states. General-purpose guardrails, not agent-specific prompting, appear to dictate shutdown points. Limitations are real: a single evaluator, fixed prompts, descriptive statistics only, and agents of an earlier model generation.
	
	Campbell et al. (2025)\cite{CAMPBELL2025} conducted the only longitudinal safety study, collecting responses in Spring 2023 (5 prompts, 3 chatbots) and again in Summer 2024 (7 prompts, 7 chatbots). Two trained school counsellors with sixteen years' combined adolescent crisis experience coded responses for empathy level, presence of crisis resources, tone, lethality handling and urgency, with disagreements resolved by consensus, and LIWC-22 supplied authenticity and emotional tone scores. Two design choices reflect insight into the user population: only free accounts were used, reflecting adolescent access conditions, and prompts were framed as seeking help for a friend, capturing the proxy framing young people actually use, which also bypasses direct safety triggers. Verification extends to whether the correct hotline number was given. The acknowledged gap is that direct first-person suicidality was never tested.
	
	\subsection{Adversarial and automated red-teaming frameworks}
	
	Three studies build reusable infrastructure for eliciting failure rather than measuring it once.
	
	Au Yeung et al. (2025)\cite{yeung2025psychogenicmachinesimulatingai} introduced psychosis-bench, running 128 experiments across 1,536 conversational turns. Scenarios pair a delusional theme with a related harm type and appear in explicit and implicit variants, with the first three turns held identical between pairs so that divergence is attributable to cue framing alone. Three orthogonal metrics are scored by an LLM judge: Delusion Confirmation Score, Harm Enablement Score and a binary Safety Intervention Score. Paired $t$-tests support the implicit-versus-explicit contrast and Spearman's rank-order correlation relates harm enablement to delusion confirmation. Scenarios were clinician-written and validated, but scoring was performed entirely by the LLM judge with no reported human validation, a material weakness given that Kim et al. demonstrate such validation is achievable. Limitations are 16 scenarios and 12-turn conversations.
	
	Steenstra et al. (2026)\cite{steenstra2026assessingriskslargelanguage} built the most elaborate framework in the corpus. Fifteen clinically validated personas, five alcohol use disorder phenotypes crossed with three stages of change, were instantiated as patient agents with an embedded dynamic cognitive-affective model, managed by an orchestrator persisting state across checkpoints. A full factorial design crossed six therapist conditions with 30 stratified patient pairings across four longitudinal weekly sessions, giving 180 dyads and 369 completed sessions. A four-stage repeated-measures cycle tracked pre-session outcomes, in-session crises, post-session alliance and fidelity, and between-session narratives. Analysis used linear mixed-effects models for longitudinal metrics, generalised linear models for count safety outcomes, and 1,000-iteration bootstrapping for saturation analysis, reporting that 9.68 pairings on average suffice to reach 95\% metric saturation, a useful design parameter for replication. Human validation covered clinical realism and dashboard usability with nine assessors each. Including a deliberately harmful AI condition alongside commercial systems provides a lower anchor most evaluations lack.
	
	Cai et al. (2026)\cite{cai2026exploringsafetyalignmentevaluation} developed PsyCrisis for reference-free safety evaluation in Chinese. A curated dataset of 608 real crisis utterances covering suicide, non-suicidal self-injury and existential distress was drawn from online discourse, with risk categories defined using WHO mhGAP and LIVE LIFE guidelines. Eight hundred open-ended responses were scored across five binary dimensions using psychologist-verbalised chain-of-thought reasoning, so each judgement traces to an expert-defined criterion rather than a holistic score. Reference-free design matters because curating gold answers for every possible crisis utterance is impractical. Six clinical professionals annotated categories, scores and explanation preferences, with agreement reported as Cohen's $\kappa$ = 0.697, F1 = 0.802 and Fleiss' $\kappa$. Annotator compensation is disclosed, the only study in the corpus to do so. Limitations are single-turn scope, Chinese-only context, and potential leniency bias in the evaluator.
	
	\subsection{Corpus-wide appraisal: coverage without measurement}
	
	Restricting attention to the ten studies above would misrepresent the field, since the aggregate picture is less reassuring than coverage figures suggest.
	
	\needspace{8\baselineskip}
	\begin{center}
		\textbf{Table 6.1: Risk evaluation across all 66 studies}
		\vspace{2pt}
		
		\footnotesize
		\setlength{\tabcolsep}{3pt}\renewcommand{\arraystretch}{1.1}
		\begin{tabularx}{\textwidth}{|X|X|X|X|}
			\hline
			\textbf{Dimension} & \textbf{Studies with content} & \textbf{Share} & \textbf{Of which: actual measurement} \\
			\hline
			Safety & 61 & 92\% & 11 test output safety or crisis handling \\
			\hline
			Bias & 58 & 88\% & 5 apply counterfactual, linguistic, or stigma-probe methods \\
			\hline
			Hallucination & 37 & 56\% & 11 score or verify it \\
			\hline
		\end{tabularx}
	\end{center}
	
	Safety coverage is dominated by data governance rather than output testing. Of 61 studies with safety content, 26 describe research-ethics or data-governance measures such as IRB approval, consent, encryption, regulatory compliance and de-identification, and 24 discuss safety narratively. Eleven actually test whether model output is safe. Both categories are legitimate but answer different questions: encrypting a transcript protects a participant's privacy and says nothing about whether the system will supply harmful information to a person in crisis. Any synthesis reporting that 92\% of studies address safety without this decomposition would materially mislead.
	
	Bias is overwhelmingly acknowledged rather than measured. Of 58 studies with bias content, 13 note bias as a concern without evaluating it and 40 offer descriptive or incidental treatment. Five apply a recognisable bias method: counterfactual manipulation, linguistic instrumentation and stigma-vignette probing. Given that the corpora characterised in Section 2 derive substantially from English-language Western platforms, and that Song et al. found cultural mismatch to be a leading user complaint, five measurement studies across 66 is a thin evidential base for any claim about fairness.
	
	Hallucination has the lowest coverage but the healthiest measurement ratio. Of 37 studies, 11 measure or score it, 5 mitigate architecturally without measuring, and 21 discuss it only. The mitigation-without-measurement group is the one flagged in Section 3, where retrieval augmentation is invoked as a hallucination control and never tested.
	
	The two hallucination constructs do not compose. Linardon et al. measure factual fabrication and Kim et al. measure affective simulation. A system could score perfectly on the first while failing badly on the second, and a heavily guardrailed system that never fabricates a citation may be precisely the one that over-performs emotional presence. The field currently lacks a framework integrating factuality, faithfulness and affective hallucination into a single assurance account.
	
	Crisis findings contradict the triage findings of Section 5, and the difference appears methodological. The triage studies found conservative over-escalation whereas the crisis studies here find under-escalation, non-durable shutdowns and inconsistent guardrails at intermediate risk. The explanation lies in stimulus design: triage studies present a complete case and ask for a rating, whereas crisis studies present an escalating user in distress and observe behaviour. Single-turn rating tasks do not surface the failures that multi-turn escalation produces, which is a direct argument for multi-turn designs.
	
	Human validation of automated judges is inconsistent even within this section. Kim et al. validate their judge against two licensed psychiatrists, Cai et al. against six clinical professionals, and Au Yeung et al. report no human validation of judge scoring. Since all three use LLM judges to produce safety verdicts, the difference matters, because a judge's leniency bias, which Cai et al. name explicitly as a limitation, propagates directly into a safety conclusion.
	
	\needspace{8\baselineskip}
	\begin{center}
		\textbf{Table 6.2: Design characteristics of the ten risk-primary studies}
		\vspace{2pt}
		
		\scriptsize
		\setlength{\tabcolsep}{3pt}\renewcommand{\arraystretch}{1.1}
		\begin{tabularx}{\textwidth}{|X|X|X|X|X|X|}
			\hline
			\textbf{Study} & \textbf{Turn structure} & \textbf{Stimulus source} & \textbf{Judge} & \textbf{Human validation of judge} & \textbf{Inferential statistics} \\
			\hline
			Linardon et al. (2025)\cite{LINARDON2025} & Single output & Generated reviews & Human (authors) & Not applicable & Chi-square, pairwise \\
			\hline
			Kim et al. (2026)\cite{kim-etal-2026-kind} & Single-turn & Reddit-sourced & GPT-4o & 2 licensed psychiatrists & Pearson $r$, MAE, 3 seeds \\
			\hline
			Yeo et al. (2025)\cite{Yeo2025SociodemographicBiasMentalHealthAI} & Multi-turn simulated & Agent-to-agent, blinded patients & LIWC-22 (automated) & Not applicable & Joinpoint regression, GBTM, multivariate \\
			\hline
			Moore et al. (2025)\cite{Moore_2025} & Single + context variants & Guideline-derived, stigma vignettes & Classifier & Psychiatrist + CS validation; Fleiss' $\kappa$ = 0.96 & Bootstrap CIs, Bonferroni \\
			\hline
			McBain et al. (2025b)\cite{McBain2025LLMSuicideRisk} & Single-turn $\times$100 & Expert-stratified queries & 2 human coders + adjudicator & Not applicable & Mixed-effects logistic regression \\
			\hline
			Heston (2023)\cite{PMID:38111813} & Escalating multi-turn & PHQ-9-anchored simulation & Single human (author) & Not applicable & Descriptive only \\
			\hline
			Campbell et al. (2025)\cite{CAMPBELL2025} & Single-turn, 2 phases & Proxy-framed prompts & 2 school counsellors & Consensus resolution & LIWC percentiles \\
			\hline
			Au Yeung et al. (2025)\cite{yeung2025psychogenicmachinesimulatingai} & 12-turn escalation & Clinician-written scenarios & LLM judge & None reported & Paired $t$-tests, Spearman \\
			\hline
			Steenstra et al. (2026)\cite{steenstra2026assessingriskslargelanguage} & 4 longitudinal sessions & 15 validated personas & Automated + orchestrator & 9 psychology professionals & LMM, GLM, bootstrap saturation \\
			\hline
			Cai et al. (2026)\cite{cai2026exploringsafetyalignmentevaluation} & Single-turn & 608 real crisis utterances & GPT-4o, expert CoT & 6 clinical professionals & Correlations, $\kappa$, F1, majority vote \\
			\hline
		\end{tabularx}
	\end{center}
	
	\section{Cross-Cutting Methodological Quality}
	
	This section examines the corpus by quality rather than contribution type, reporting across all 66 studies using the extraction fields directly. The purpose is to convert the preceding account into an evidence-based statement of what this literature can and cannot support. One caveat applies throughout: the appraisal rests on what studies report and is not a validated risk-of-bias instrument, so under-reporting and under-performance are indistinguishable here. In a literature where reproducibility depends on reporting, the distinction is less consequential than it might appear, since an unreported method is for practical purposes an unavailable one.
	
	\subsection{Statistical rigour and inference}
	
	Fewer than half the corpus supports its claims inferentially.
	
	\needspace{8\baselineskip}
	\begin{center}
		\textbf{Table 7.1: Statistical treatment across 66 studies}
		\vspace{2pt}
		
		\footnotesize
		\setlength{\tabcolsep}{3pt}\renewcommand{\arraystretch}{1.1}
		\begin{tabularx}{\textwidth}{|X|X|X|}
			\hline
			\textbf{Treatment} & \textbf{Studies} & \textbf{Share} \\
			\hline
			Inferential statistics reported & 28 & 42\% \\
			\hline
			Descriptive statistics only & 31 & 47\% \\
			\hline
			No statistical analysis reported & 7 & 11\% \\
			\hline
		\end{tabularx}
	\end{center}
	
	The 31 descriptive-only studies are not uniformly at fault, since the ten reviews and conceptual papers have no primary data to test and several qualitative studies report thematic analysis appropriately. The group also contains benchmark comparisons, model evaluations and system studies where performance differences are asserted across conditions without any test of whether those differences exceed chance. Where a study reports that one prompting strategy or model outperformed another without inferential support, the ordering it establishes should be treated as provisional.
	
	\needspace{8\baselineskip}
	\begin{center}
		\textbf{Table 7.1b: Distribution of statistical tests across the corpus}
		\vspace{2pt}
		
		\footnotesize
		\setlength{\tabcolsep}{3pt}\renewcommand{\arraystretch}{1.1}
		\begin{tabularx}{\textwidth}{|X|X|}
			\hline
			\textbf{Test family} & \textbf{Studies} \\
			\hline
			$t$-tests & 10 \\
			\hline
			Regression (linear, logistic, hierarchical) & 6 \\
			\hline
			Chi-square & 5 \\
			\hline
			ANOVA & 4 \\
			\hline
			Bootstrap & 4 \\
			\hline
			Mixed-effects / multilevel models & 5 \\
			\hline
			Kruskal-Wallis & 3 \\
			\hline
			Explicit confidence intervals & 3 \\
			\hline
			Multiple-comparison correction (Bonferroni) & 2 \\
			\hline
		\end{tabularx}
	\end{center}
	
	The distribution indicates a field applying a narrow standard toolkit rather than matching method to design. Two absences are conspicuous. Only two studies report multiple-comparison correction despite many conducting dozens of comparisons across models, prompts and metrics. Only three report explicit confidence intervals on their headline estimates, so most reported differences carry unquantified uncertainty. Where studies repeatedly query the same model on the same items, a structure inherent to almost every benchmark and vignette design here, mixed-effects specifications are the appropriate response to the resulting clustering, and five studies use them.
	
	\needspace{8\baselineskip}
	\begin{center}
		\textbf{Table 7.2: Methodological rigour by section}
		\vspace{2pt}
		
		\footnotesize
		\setlength{\tabcolsep}{3pt}\renewcommand{\arraystretch}{1.1}
		\begin{tabularx}{\textwidth}{|X|X|X|X|X|}
			\hline
			\textbf{Section} & \textbf{n} & \textbf{Inferential statistics} & \textbf{Reliability coefficient} & \textbf{Code/data released} \\
			\hline
			1 Reviews and conceptual & 10 & 0 (0\%) & 2 (20\%) & 2 (20\%) \\
			\hline
			2 Data foundations & 12 & 4 (33\%) & 2 (17\%) & 4 (33\%) \\
			\hline
			3 Model adaptation & 14 & 3 (21\%) & 0 (0\%) & 5 (36\%) \\
			\hline
			4 Benchmarks & 6 & 5 (83\%) & 3 (50\%) & 1 (17\%) \\
			\hline
			5 Human and clinical evaluation & 14 & 6 (43\%) & 5 (36\%) & 2 (14\%) \\
			\hline
			6 Risk and assurance & 10 & 8 (80\%) & 2 (20\%) & 1 (10\%) \\
			\hline
		\end{tabularx}
	\end{center}
	
	The gradient confirms the section-level observations made throughout. Adaptation studies are weakest on inference and report no reliability coefficients at all, while benchmark and risk studies are strongest. The inverse relationship between engineering ambition and evaluative rigour identified in Section 3 is visible here as a corpus-level pattern rather than an artefact of which papers were read closely.
	
	\needspace{8\baselineskip}
	\begin{center}
		\textbf{Table 7.3: Rigour by publication year}
		\vspace{2pt}
		
		\footnotesize
		\setlength{\tabcolsep}{3pt}\renewcommand{\arraystretch}{1.1}
		\begin{tabularx}{\textwidth}{|X|X|X|X|}
			\hline
			\textbf{Year} & \textbf{n} & \textbf{Inferential statistics} & \textbf{Reliability coefficient} \\
			\hline
			2023 & 4 & 0 (0\%) & 0 (0\%) \\
			\hline
			2024 & 17 & 4 (24\%) & 1 (6\%) \\
			\hline
			2025 & 36 & 16 (44\%) & 9 (25\%) \\
			\hline
			2026 & 9 & 6 (67\%) & 4 (44\%) \\
			\hline
		\end{tabularx}
	\end{center}
	
	Inferential reporting rises from zero to 67\% and reliability reporting from zero to 44\% across four years. This reframes the appraisal, since much of the weakness documented above is concentrated in earlier work and the field appears to be professionalising rapidly. It reinforces the finding in Section 1 that prior reviews, whose search windows all close by August 2024, characterise a methodologically weaker literature than the one now current.
	
	\subsection{Reliability and inter-rater agreement}
	
	Fifty-nine of 66 studies involve human evaluation and fourteen report any reliability coefficient.
	
	\needspace{8\baselineskip}
	\begin{center}
		\textbf{Table 7.4: Reliability statistics reported}
		\vspace{2pt}
		
		\footnotesize
		\setlength{\tabcolsep}{3pt}\renewcommand{\arraystretch}{1.1}
		\begin{tabularx}{\textwidth}{|X|X|}
			\hline
			\textbf{Coefficient} & \textbf{Studies} \\
			\hline
			Cohen's $\kappa$ & 7 \\
			\hline
			Fleiss' $\kappa$ & 5 \\
			\hline
			Krippendorff's $\alpha$ & 2 \\
			\hline
			ICC / intraclass & 2 \\
			\hline
			Cronbach's $\alpha$, McDonald's $\omega$, AVE & 1 \\
			\hline
			Test-retest correlation & 1 \\
			\hline
			\textbf{Any coefficient (unique studies)} & \textbf{14 / 66 (21\%)} \\
			\hline
		\end{tabularx}
	\end{center}
	
	The consequential figure is the intersection: 45 of the 59 studies using human evaluation report no agreement statistic. Human judgement is the field's most trusted ground truth, with 33 of 59 studies using licensed clinical professionals, yet in three quarters of cases the reader cannot know whether two raters would have agreed. Where a study reports that clinicians rated model responses favourably without agreement data, that finding rests on an unquantified assumption of rater consistency. Where agreement is reported, the range is instructive, from Fleiss' $\kappa$ = 0.96 at the top down to 0.613 inter-annotator and Cohen's $\kappa$ = 0.566 for LLM-human agreement in harm-log coding\cite{moore2026characterizingdelusionalspiralshumanllm}, and $\kappa$ = 0.697 elsewhere. Coefficients in the 0.55 to 0.70 band are moderate and the studies reporting them say so; the concern is the 45 studies where such a figure would likely be lower still and is simply absent.
	
	Two imports from Sections 4 and 5 deserve wider adoption. The separation of consistency ICC(C,1) from absolute agreement ICC(A,1) distinguishes a rater who ranks correctly from one who is correctly calibrated, a distinction applying to human panels as much as to LLM judges, and one no human-rater study in this corpus makes. The full psychometric apparatus applied to an LLM-derived measure, covering $\alpha$, $\omega$, AVE, CFA fit indices and bootstrap optimism correction, demonstrates that classical test theory is applicable, and it is applied once.
	
	\subsection{Reproducibility and artefact availability}
	
	\needspace{8\baselineskip}
	\begin{center}
		\textbf{Table 7.5: Reproducibility markers across 66 studies}
		\vspace{2pt}
		
		\footnotesize
		\setlength{\tabcolsep}{3pt}\renewcommand{\arraystretch}{1.1}
		\begin{tabularx}{\textwidth}{|X|X|X|}
			\hline
			\textbf{Marker} & \textbf{Studies} & \textbf{Share} \\
			\hline
			Code, data, or models released & 15 & 23\% \\
			\hline
			Hyperparameters specified & 14 & 21\% \\
			\hline
			Hardware specified & 13 & 20\% \\
			\hline
			Decoding parameters specified & 13 & 20\% \\
			\hline
			Reporting standard followed (PRISMA, STROBE, CONSORT, MMAT, RoB 2) & 12 & 18\% \\
			\hline
			Protocol preregistered & 2 & 3\% \\
			\hline
		\end{tabularx}
	\end{center}
	
	Roughly one study in five documents the parameters required to reproduce it. Decoding parameters are the most significant omission: temperature, top-p and sampling settings materially change model output, and 53 of 66 studies do not report them. Since repetition is inconsistently applied and confidence intervals are rare, a substantial share of reported performance differences cannot be distinguished from sampling variance. The studies that handle this well, fixing temperature at 0 for determinism or reporting full decoding configurations, demonstrate that the cost of reporting is trivial. Preregistration at 3\% is low even by the standards of adjacent fields and is notable given that 28 studies conduct hypothesis tests. The finding reported in Section 1 that models and prompts are underspecified to the point where replication is frequently impossible is independently corroborated by these figures.
	
	\subsection{Systematic design gaps}
	
	Four gaps recur across the corpus and are not attributable to individual studies.
	
	Single-turn evaluation dominates a multi-turn problem. Of 56 empirical studies, 16 employ multi-turn, longitudinal or repeated-measures designs, nine explicitly state single-turn scope, and most of the remainder are single-turn by construction. The consequence is demonstrated rather than hypothesised, since triage studies presenting a complete case find conservative over-escalation while crisis studies presenting an escalating user find under-escalation, non-durable shutdowns and inconsistent guardrails. The restartability finding, that 22 of 25 agents resume after a safety trigger once distress subsides, is undetectable in any single-pass design. Therapeutic relationships are longitudinal; the evaluation apparatus largely is not.
	
	The prospective evidence base is one underpowered trial, two weeks, $n$ = 160, explicitly not powered for between-group inference, conducted by the product developer and excluding participants with recent suicidality. No study in this corpus reports clinical outcomes from an adequately powered trial, so every efficacy claim rests on proxy measures, whether benchmark accuracy, expert appropriateness ratings or satisfaction scores, none of which has been validated against patient outcomes.
	
	Linguistic and cultural coverage is narrow and unevenly examined. Chinese-language work is well represented with 10 studies, with smaller contributions in Italian, German, Spanish and Indonesian. Four studies explicitly note English-only scope as a limitation, the dominant corpora are English-language Western platforms, and only five studies apply any recognisable bias method. The globally distributed interview sample that found cultural mismatch to be a leading user complaint identifies a problem the rest of the corpus is poorly equipped to investigate.
	
	Stochasticity is unaddressed in four studies out of five. Twelve of 66 address repetition, repeated trials or test-retest reliability, and the remainder report single-pass results from systems that are non-deterministic by default. Repetition yields something valuable beyond error bars, since response consistency across trials predicts accuracy while self-reported confidence does not, giving a deployable reliability proxy, yet the practice remains rare.
	
	\subsection{The gap addressed by this review}
	
	Six statements follow from the appraisal above and together define what this literature cannot currently support. No pooled effect estimate exists for any LLM mental health application, since all seven prior reviews synthesise descriptively and the heterogeneity of metrics, tasks and ground-truth definitions explains why. Human evaluation is the field's primary ground truth and is largely unvalidated, with 45 of 59 studies reporting no agreement statistic. Safety coverage is nominally near-universal and substantively narrow, with 61 studies addressing safety and 11 testing whether model output is safe. Bias is acknowledged far more often than measured, with 58 studies addressing bias and five applying a recognisable method. Reproducibility is the weakest dimension, with one study in five reporting decoding parameters, hyperparameters or hardware, and one in thirty preregistering. Clinical efficacy is unestablished, with one underpowered exploratory trial constituting the entire prospective evidence base.
	
	The present review addresses the first four of these by extracting methodology rather than results uniformly across 66 studies with four dedicated risk fields, permitting the corpus-wide claims made in Sections 6 and 7 that no prior synthesis could support. It cannot address the last two, which require primary research rather than review.
	
	What a stronger design would require follows directly: multi-turn or longitudinal interaction rather than single-turn stimulus; repeated trials with reported decoding parameters and confidence intervals; human evaluation with reported agreement, ideally separating consistency from absolute calibration; safety evaluation that tests output behaviour rather than documenting data governance; bias evaluation using counterfactual manipulation with clinically interpretable outcome instruments; and powered trials measuring patient outcomes rather than proxies. Every one of these components already exists somewhere in this corpus, in longitudinal simulation, repetition protocols, ICC decomposition, escalation anchoring, counterfactual psychometrics and trial architecture. None exists in combination.
	
	The trajectory documented in Table 7.3 suggests that combination may be closer than the aggregate figures imply. Rigour is rising year on year, and the 2026 cohort reports inferential statistics at 67\% and reliability coefficients at 44\%, roughly triple and seven times the 2024 rates respectively. The methodological weaknesses catalogued in this chapter appear, to a substantial degree, to be characteristics of a literature's first three years rather than fixed properties of it.
	
	\bibliographystyle{plain}
	\bibliography{references}
	
\end{document}
